\section{Relative ample Nadel vanishing and the smooth envelope theorem in characteristic zero}
\label{sec:kn-relative-nadel-vanishing}

This section combines the relative Nadel assembly
(\cref{thm:kn-relative-ample-nadel-of-cover-trace}), whose only hypothesis is the
canonical trace input on adapted covers, with the existence of that canonical
trace in characteristic zero (\cref{thm:kn-cover-canonical-trace-input}). The
result is relative ample Nadel vanishing over every finite-type affine base,
the statement of \cref{def:c0n-relative-ample-nadel-input}. Through the
multiplier-ideal argument of \cref{thm:c0n-original-smooth-envelope-from-nadel}
it gives the envelope property of smooth projectively normal subvarieties of
$\mathbf P^M$ in characteristic zero, the theorem that \cite[\S3]{FHX} cites
from \cite[Theorem~5.20]{BJ22}, and therefore that envelope property over every
algebraically closed field.

\begin{theorem}[Relative ample Nadel vanishing in characteristic zero {\cite{EV92}}]
\label{thm:kn-relative-ample-nadel-vanishing}
\uses{thm:kn-relative-ample-nadel-of-cover-trace,thm:kn-cover-canonical-trace-input,def:c0n-relative-ample-nadel-input}
\lean{EnvelopeProperty.Char0.relativeAmpleNadelVanishingInput}
\leanok
Let $k$ be a field of characteristic zero and $R$ a finite-type $k$-algebra.
Let $X$ be an integral scheme, smooth of relative dimension $n$ over $k$, with
$\pi:X\to\operatorname{Spec}R$ projective and compatible with the structure maps
to $\operatorname{Spec}k$. Let $I\subseteq\mathcal O_X$ be an ideal that is
nonzero at the generic point, let $A,B$ be invertible with $I\otimes B$
generated by finitely many global sections, let $c\geq0$ be rational, and let
$d>0$, $t\geq0$ be integers with $dc=t$ and
$A^{\otimes d}\simeq C\otimes B^{\otimes t}$ for a line $C$ that is ample over
$\operatorname{Spec}R$. Then
\[
 H^{q}\bigl(X,\;(\omega_{X/k}\otimes A)\otimes\mathcal J(c\cdot I)\bigr)=0
 \qquad(q>0),
\]
where $\mathcal J(c\cdot I)$ is the multiplier ideal defined by all divisorial
valuations. In other words, the relative ample Nadel input
(\cref{def:c0n-relative-ample-nadel-input}) holds for every field of
characteristic zero.
\end{theorem}
\begin{proof}
By \cref{thm:kn-cover-canonical-trace-input} the canonical trace input
(\cref{def:kn-relative-nadel-cover-trace-input}) holds over $k$, since $k$ has
characteristic zero. By \cref{thm:kn-relative-ample-nadel-of-cover-trace} the
canonical trace input over $k$ implies the relative ample Nadel input over $k$,
which is the displayed vanishing for all data as in the statement.
\end{proof}

\begin{theorem}[Envelopes of smooth projectively normal subvarieties in characteristic zero {\cite[Theorem~5.20]{BJ22}, \cite[Theorem~3.1]{FHX}}]
\label{thm:kn-smooth-envelope-char-zero}
\uses{thm:c0n-original-smooth-envelope-from-nadel,thm:kn-relative-ample-nadel-vanishing,conv:sg-setting,def:sg-ep}
\lean{EnvelopeProperty.smoothEnvelopeCharZero_unconditional,EnvelopeProperty.SmoothEnvelopeCharZero.instOfIsAlgClosed}
\leanok
Let $k$ be an algebraically closed field of characteristic zero and let
$X=V(I)\subseteq\mathbf P^M_k$ be projectively normal with smooth charts. Then
every class $\theta\,c_1(\mathcal O_X(1))$ with $\theta\geq0$ has the envelope
property on $X$.
\end{theorem}
\begin{proof}
By \cref{thm:kn-relative-ample-nadel-vanishing} the relative ample Nadel input
holds over $k$. Theorem~\ref{thm:c0n-original-smooth-envelope-from-nadel} derives from
that input, over an algebraically closed field of characteristic zero, the
envelope property of every projectively normal $X=V(I)\subseteq\mathbf P^M_k$
with smooth charts in every class $\theta\,c_1(\mathcal O_X(1))$,
$\theta\geq0$.
\end{proof}

\begin{corollary}[Envelopes of smooth projectively normal subvarieties in every characteristic {\cite[Theorem~3.1]{FHX}}]
\label{cor:kn-smooth-envelope-all-characteristics}
\uses{thm:sg-ep,thm:kn-smooth-envelope-char-zero}
\lean{EnvelopeProperty.ProjSubvariety.hasEnvelopeProperty_of_smooth_unconditional}
\leanok
Let $k$ be an algebraically closed field and let $X=V(I)\subseteq\mathbf P^M_k$
be projectively normal with smooth charts. Then every class
$\theta\,c_1(\mathcal O_X(1))$, $\theta\in\mathbf R$, has the envelope property
on $X$.
\end{corollary}
\begin{proof}
Theorem~\ref{thm:sg-ep} proves the statement in characteristic $p>0$, and in
characteristic zero under the characteristic-zero envelope statement for
$\theta\geq0$; the latter is \cref{thm:kn-smooth-envelope-char-zero}.
\end{proof}
